% Generated draft: requires mathematical and source review.
\documentclass[11pt]{article}
\usepackage[margin=25mm]{geometry}
\usepackage{fontspec}
\setmainfont{Latin Modern Roman}
\usepackage{amsmath,amssymb,mathtools,unicode-math}
\setmathfont{Latin Modern Math}
\usepackage{enumitem,needspace,xurl,hyperref}
\setlength{\parindent}{0pt}
\setlength{\parskip}{5pt}
\setlength{\emergencystretch}{3em}
\newcommand{\N}{\mathbb{N}}
\newcommand{\Z}{\mathbb{Z}}
\newcommand{\Q}{\mathbb{Q}}
\newcommand{\R}{\mathbb{R}}
\newcommand{\C}{\mathbb{C}}
\newcommand{\F}{\mathbb{F}}
\newcommand{\A}{\mathbb{A}}
\newcommand{\PP}{\mathbb{P}}
\newcommand{\E}{\mathbb{E}}
\newcommand{\eps}{\varepsilon}
\newcommand{\dd}{\,\mathrm{d}}
\begin{document}
\section*{UnsolvedMath ID: 6200073}
\textbf{Source problem code:} catalogue entry
\subsection{Definitions and mathematical statement}
For rank \(n\ge2\), projectivized Culler--Vogtmann Outer space \(CV_n\) is the space of finite marked metric graphs of rank \(n\), with positive edge lengths normalized to total length one, modulo marking-preserving isometry. A marking is a homotopy equivalence from the rank-\(n\) rose; standard graph reductions remove inessential vertices. The source does not specify rank or whether this normalized or unprojectivized version is intended.

A Poisson boundary belongs to a Markov process, not to a topological space alone. For a chosen kernel \(P\) on a state space \(X\), bounded harmonic functions satisfy
\[
h(x)=\int_Xh(y)\,P(x,\mathrm dy).
\]
A boundary model \((B,\nu_x)\) represents these functions by the Poisson transform
\[
h(x)=\int_B F(\xi)\,\nu_x(\mathrm d\xi),
\]
with the appropriate uniqueness/maximality property for bounded boundary functions. Identifying this measured harmonic boundary, with its hitting measures, is different from naming a topological compactification.

Duchin's question asks for the Poisson boundary of Outer space, but gives no random walk, diffusion, transition measure, or starting-law convention. If the intended process is induced by a random walk on \(\operatorname{Out}(F_n)\) acting on \(CV_n\), its increment law and hypotheses, and the distinction between the group walk and its orbit process, must be fixed. A group Poisson-boundary theorem is not automatically the answer for every process on Outer space.

The supplied research summary reports partial progress involving Horbez and free-factor/sphere-complex-related boundary models. Its description of a Poisson boundary as a flag complex does not specify a measured boundary or transition law and is not adopted as a precise identification. No exact process-specific conclusion can be extracted from that discussion.
\subsection{Short English statement}
Identify the measured Poisson boundary of an intended random process on Outer space. The source supplies no process or probability law, so Outer space alone does not determine the answer.
\subsection{Sources}
\begin{itemize}
\item Local UnsolvedMath registry, permanent ID \(6200073\), exact source code \texttt{catalogue entry}, ``Boundaries of Groups and Kleinian Groups --- Problem 73.'' Proposer Moon Duchin, 2005.
\item Kapovich, \emph{Problems on Boundaries of Groups and Kleinian Groups} (2005), Problem 73, recorded PDF page 19; \url{https://www.math.ucdavis.edu/~kapovich/EPR/problems.pdf}. Recorded access 2026-07-29; pdf-text-normalized extraction.
\item Supplied PARTIAL-PROGRESS research summary names Horbez and discusses free-group, \(\operatorname{Out}(F_n)\), and free-factor/sphere-complex boundaries, but supplies no complete publication citation, increment law, or theorem hypotheses.
\end{itemize}
\end{document}
